\documentclass[11pt,letterpaper]{article}
\usepackage[margin=0.72in]{geometry}
\usepackage[T1]{fontenc}
\usepackage[utf8]{inputenc}
\usepackage{lmodern}
\usepackage{microtype}
\usepackage{xcolor}
\usepackage{enumitem}
\usepackage{tabularx}
\usepackage{hyperref}

\definecolor{accent}{HTML}{00419B}
\definecolor{muted}{HTML}{555555}

\hypersetup{
  colorlinks=true,
  urlcolor=accent,
  linkcolor=accent
}

\setlength{\parindent}{0pt}
\setlength{\parskip}{0.45em}
\setlist[itemize]{leftmargin=1.25em, topsep=0.2em, itemsep=0.15em}

\newcommand{\MauriName}{Maurizio Ungaro}
\newcommand{\MauriTitle}{Staff Scientist, Jefferson Lab}
\newcommand{\MauriEmail}{\href{mailto:ungaro@jlab.org}{ungaro@jlab.org}}
\newcommand{\MauriHomepage}{\href{https://maureeungaro.github.io/home/}{maureeungaro.github.io/home}}
\newcommand{\MauriGitHub}{\href{https://github.com/maureeungaro}{github.com/maureeungaro}}
\newcommand{\MauriScholar}{\href{https://scholar.google.com/citations?user=zkWYILYAAAAJ\&hl=en}{Google Scholar}}
\newcommand{\MauriInspire}{\href{https://inspirehep.net/authors/1322331}{INSPIRE}}
\newcommand{\MauriOrcid}{\href{https://orcid.org/0000-0001-6982-3310}{ORCID 0000-0001-6982-3310}}
\newcommand{\MauriLinkedIn}{\href{https://www.linkedin.com/in/maurizio-ungaro-37992062/}{LinkedIn}}
\newcommand{\MauriJLab}{\href{https://www.jlab.org}{Jefferson Lab}}
\newcommand{\MauriHallB}{\href{https://www.jlab.org/physics/hall-b}{Jefferson Lab Hall-B}}
\newcommand{\MauriRPI}{\href{https://www.rpi.edu}{Rensselaer Polytechnic Institute}}
\newcommand{\MauriUniGe}{\href{https://unige.it/}{Universita degli Studi di Genova}}
\newcommand{\MauriGEMC}{\href{https://gemc.github.io/home/}{GEMC}}
\newcommand{\MauriGeantFourJLab}{\href{https://jeffersonlab.github.io/g4home/}{Geant4 at JLab}}
\newcommand{\MauriClasTwelve}{\href{https://github.com/gemc/clas12Tags}{CLAS12 Simulations}}
\newcommand{\MauriOSG}{\href{https://maureeungaro.github.io/home/osg/osg}{CLAS12 on OSG}}
\newcommand{\MauriLTCC}{\href{https://maureeungaro.github.io/home/ltcc/ltcc}{Low Threshold Cherenkov Counter}}
\newcommand{\MauriSPringEight}{\href{https://www.spring8.or.jp/en/}{SPring-8}}
\newcommand{\MauriHPS}{\href{https://www.jlab.org/physics/hall-b/other}{Heavy Photon Search}}

\newcommand{\DocHeader}[1]{
  {\LARGE\bfseries \MauriName}\hfill {\color{muted}#1}\\
  Staff Scientist, \MauriJLab{}\\
  \MauriEmail{} \quad
  \MauriHomepage{} \quad
  \MauriGitHub{}\\
  \MauriScholar{} \quad
  \MauriInspire{} \quad
  \MauriOrcid{}\\
  \vspace{0.4em}
  \hrule
  \vspace{0.8em}
}

\newcommand{\RoleEntry}[4]{
  \textbf{#1}\hfill {\color{muted}#2}\\
  #3\\
  #4
}

\newcommand{\ProjectEntry}[3]{
  \textbf{#1}: #2\\
  {\small\href{#3}{#3}}
}


\begin{document}
\DocHeader{Scientific CV}

\section*{Profile}
Nuclear physicist and simulation software developer at Jefferson Lab working on
Geant4-based detector simulations, \MauriGEMC{}, CLAS12 production workflows, detector
systems, and nucleon-structure research. My work connects physics analysis,
detector operations, and simulation infrastructure, with a focus on making
Geant4 workflows easier to build, run, document, and share.

\section*{Current Position}
\RoleEntry{Staff Scientist}{2011--present}{\MauriJLab{}, Experimental Hall-B, Newport News, VA}
{Hall-B and CLAS12 detector operations, simulation software, Geant4 support, LTCC calibration and performance, distributed simulation production, and nucleon-structure analysis.}

\section*{Previous Positions}
\RoleEntry{Postdoctoral Research Associate and Research Associate}{2004--2011}{University of Connecticut, Storrs, CT}
{CLAS analyses, meson electro-production, detector software, and CLAS12 simulation and reconstruction development.}

\section*{Education}
\begin{tabularx}{\textwidth}{@{}lX@{}}
2003 & Ph.D. in Nuclear Physics, \MauriRPI{}, Troy, NY. Thesis: pi0 electro-production from Delta(1232) at high momentum transfer.\\
1999 & Laurea in Fisica, \MauriUniGe{}, Genoa, Italy.\\
\end{tabularx}

\section*{Research and Technical Leadership}
\begin{itemize}
  \item Nucleon-structure measurements using meson electro-production and CLAS/CLAS12 data.
  \item Development and support of \MauriGEMC{}, a database-driven Geant4 simulation workflow.
  \item CLAS12 simulation releases, geometry tags, and production workflows through \MauriClasTwelve{}.
  \item Distributed simulation production using HTCondor and Open Science Grid resources.
  \item CLAS12 \MauriLTCC{} operation, maintenance, calibration, and detector-performance studies.
  \item \MauriGeantFourJLab{} tutorials, examples, and user support for Jefferson Lab users.
\end{itemize}

\section*{Selected Projects}
\ProjectEntry{\MauriGEMC{}}{Database-driven Geant4 simulations with a Python-friendly workflow.}{https://gemc.github.io/home/}

\ProjectEntry{\MauriGeantFourJLab{}}{Tutorials, examples, and support material for Jefferson Lab users.}{https://jeffersonlab.github.io/g4home/}

\ProjectEntry{\MauriClasTwelve{}}{Tagged CLAS12 geometry and simulation releases.}{https://github.com/gemc/clas12Tags}

\ProjectEntry{\MauriOSG{}}{Distributed CLAS12 simulation production workflows.}{https://maureeungaro.github.io/home/osg/osg}

\ProjectEntry{\MauriLTCC{}}{CLAS12 forward detector work for pion/kaon discrimination.}{https://maureeungaro.github.io/home/ltcc/ltcc}

\section*{Selected Collaborations and Facilities}
\begin{itemize}
  \item \MauriJLab{} and \MauriHallB{}: current laboratory and experimental program.
  \item \MauriHPS{}: dark-sector search collaboration at Jefferson Lab.
  \item \MauriSPringEight{}: prior experimental work in Japan through LEPS/SPring-8 collaboration.
\end{itemize}

\section*{Technical Skills}
\textbf{Simulation and analysis}: Geant4, GEMC, CLAS12 simulations, ROOT, detector geometry, digitization workflows, PyVista visualization.

\textbf{Programming and infrastructure}: C++, Python, shell scripting, Git, GitHub, CI, Docker, Meson, CMake, SCons.

\textbf{Scientific computing}: HTCondor, Open Science Grid, SQLite, CSV/JSON/ROOT workflows, reproducible environments, documentation.

\textbf{Communication}: tutorials, technical writing, presentations, mentoring.

\section*{Selected Publications}
\begin{itemize}
  \item \textit{Geant4 Monte-Carlo (GEMC): A database-driven simulation program}, EPJ Web of Conferences 295, 05005 (2024).
  \item \textit{The CLAS12 Geant4 simulation}, Nucl. Instrum. Meth. A 959 (2020).
  \item \textit{The CLAS12 Low Threshold Cherenkov detector}, Nucl. Instrum. Meth. A 957 (2020).
  \item \textit{Measurement of the N to Delta(1232) transition at high momentum transfer by pi0 electro-production}, Phys. Rev. Lett. 97 (2006).
\end{itemize}

\section*{Online Profiles}
\MauriHomepage{} \quad
\MauriGitHub{} \quad
\MauriScholar{} \quad
\MauriInspire{} \quad
\MauriOrcid{}

\end{document}
